% Generated draft: requires mathematical and source review.
\documentclass[11pt]{article}
\usepackage[margin=25mm]{geometry}
\usepackage{fontspec}
\setmainfont{Latin Modern Roman}
\usepackage{amsmath,amssymb,mathtools,unicode-math}
\setmathfont{Latin Modern Math}
\usepackage{enumitem,needspace,xurl,hyperref}
\setlength{\parindent}{0pt}
\setlength{\parskip}{5pt}
\setlength{\emergencystretch}{3em}
\newcommand{\N}{\mathbb{N}}
\newcommand{\Z}{\mathbb{Z}}
\newcommand{\Q}{\mathbb{Q}}
\newcommand{\R}{\mathbb{R}}
\newcommand{\C}{\mathbb{C}}
\newcommand{\F}{\mathbb{F}}
\newcommand{\A}{\mathbb{A}}
\newcommand{\PP}{\mathbb{P}}
\newcommand{\E}{\mathbb{E}}
\newcommand{\eps}{\varepsilon}
\newcommand{\dd}{\,\mathrm{d}}
\begin{document}
\section*{UnsolvedMath ID: 800666}
\textbf{Source problem code:} catalogue entry
\subsection{Definitions and mathematical statement}
The supplied record asks about a polarized partition relation for open \(\omega\)-covers. Its excerpt does not specify the ambient class of spaces, the convention on countability of covers, or the quantification of \(k\). The following gives a precise parameterized formulation under the Cartesian-product interpretation of the displayed relation; these missing conventions require confirmation from the source.

Let \(\mathscr X\) be a specified class of topological spaces. For each \(X\in\mathscr X\), let \(\mathscr U_X\) specify the admissible open covers of \(X\). Define
\[
\Omega(X)=\left\{\mathcal U\in\mathscr U_X:X\notin\mathcal U\ \text{and}\ (\forall F\subseteq X\text{ finite})(\exists U\in\mathcal U)\ F\subseteq U\right\}.
\]
Thus every finite subset of \(X\) is contained in one member of an \(\omega\)-cover, and no member is the entire space. Whether \(\mathscr U_X\) consists of all open covers or only countable open covers is left unresolved by the supplied excerpt.

The property \(\mathsf S_1(\Omega,\Omega)\) means
\[
(\forall (\mathcal U_n)_{n\geq1}\in\Omega(X)^{\{1,2,\ldots\}})
(\exists (U_n)_{n\geq1})\quad
(\forall n\geq1)\ U_n\in\mathcal U_n
\quad\text{and}\quad
\{U_n:n\geq1\}\in\Omega(X).
\]
For a positive integer \(k\), write \(P_k(X)\) for the following Cartesian-product reading of
\[
\binom{\Omega}{\Omega}\longrightarrow
\left[\begin{matrix}\Omega\\\Omega\end{matrix}\right]_{k/<3}:
\]
for every \(\mathcal U,\mathcal V\in\Omega(X)\) and every coloring
\(c:\mathcal U\times\mathcal V\to\{1,\ldots,k\}\), there are subfamilies
\(\mathcal U'\subseteq\mathcal U\) and \(\mathcal V'\subseteq\mathcal V\) such that
\[
\mathcal U',\mathcal V'\in\Omega(X)
\quad\text{and}\quad
\bigl|c[\mathcal U'\times\mathcal V']\bigr|<3.
\]
In particular, the retained rectangle uses at most two colors. This definition treats the coloring as defined on the entire ordered Cartesian product; the excerpt supplies no definition confirming this convention or specifying additional restrictions.

Let \(P\) denote the property intended by the source notation, once its conventions have been fixed. The first question is whether
\[
(\forall X\in\mathscr X)\qquad
P(X)\ \Longleftrightarrow\ X\models\mathsf S_1(\Omega,\Omega).
\]
The excerpt does not determine whether \(P=P_k\) for a fixed \(k\), or whether \(P(X)\) means \((\forall k\geq1)P_k(X)\). For \(k=1,2\), the displayed Cartesian-product condition is automatic, so this distinction is material.

For a property \(Q\), define its critical cardinality relative to \(\mathscr X\), when the minimum exists, by
\[
\operatorname{non}_{\mathscr X}(Q)
=\min\{|X|:X\in\mathscr X\text{ and }X\text{ fails }Q\}.
\]
Let \(\mathcal M\) be the ideal of meager subsets of \(\R\), and define
\[
\mathsf{cov}(\mathcal M)
=\min\left\{|\mathcal A|:\mathcal A\subseteq\mathcal M,
\ \bigcup\mathcal A=\R\right\}.
\]
If the proposed equivalence fails, the second question is whether
\[
\operatorname{non}_{\mathscr X}(P)
=\operatorname{non}_{\mathscr X}\bigl(\mathsf S_1(\Omega,\Omega)\bigr)
=\mathsf{cov}(\mathcal M).
\]
The identification of the critical cardinality of \(\mathsf S_1(\Omega,\Omega)\) with \(\mathsf{cov}(\mathcal M)\) is supplied parenthetically by the source. It is recorded as source context, subject to the source's missing ambient conventions, rather than asserted here for an arbitrary \(\mathscr X\).
\subsection{Short English statement}
The source asks whether a polarized coloring property for open \(\omega\)-covers is equivalent to selecting one member from each of a sequence of \(\omega\)-covers so that the selected family is again an \(\omega\)-cover. If equivalence fails, it asks whether the least size of a counterexample is nevertheless \(\mathsf{cov}(\mathcal M)\). The ambient spaces, cover convention, and quantification of \(k\) require clarification.
\subsection{Sources}
\begin{itemize}
\item Local UnsolvedMath registry, permanent ID \(800666\), exact source problem code \texttt{catalogue entry}; collection: AMR Open Problem Lists. The supplied record labels the problem \texttt{open}; this is a supplied status, not an independently verified current assessment.
\item Tsaban, \emph{Selection Principles in Mathematics: A Milestone of Open Problems} (2003), Problem 39. Supplied source URL: \url{https://arxiv.org/abs/math/0312182}. The record reports access on 2026-10-16 and extraction by \texttt{arxiv\_latex}; these are supplied provenance details, not a claim of access or verification for this draft.
\item The source gives the attribution/context label \texttt{[PolarPar, Problem 1]}. No full bibliographic entry for \texttt{PolarPar} is included in the supplied record. The registry attributes the problem to Tsaban and records the year 2003; the underlying attribution cannot be resolved from the excerpt alone.
\end{itemize}
\end{document}
